%======================================================================
% فصل ششم: نتیجه‌گیری
%======================================================================
\chapter{نتیجه‌گیری}
\label{ch:conclusion}

\section{جمع‌بندی}

این پایان‌نامه به مسئله‌ی یکپارچه‌سازی سرویس‌های هوش مصنوعی از منظر
فنی، مالی و مدیریتی پرداخت. مسئله‌ی اصلی آن بود که سازمان‌ها و
توسعه‌دهندگان، برای استفاده از چند ارائه‌دهنده‌ی هوش مصنوعی، ناچار به
پذیرش چندگانگی در \lr{API}، مدل مالی، گزارش‌گیری و کنترل دسترسی هستند.
برای پاسخ به این مسئله، یک سامانه‌ی واسط مبتنی بر \lr{Laravel} طراحی و
پیاده‌سازی شد که رابطی سازگار با \lr{OpenAI} را در سطح بیرونی ارائه می‌دهد
و در سطح درونی، قابلیت‌های احراز هویت، مدیریت کلید، کیف پول، اشتراک،
فاکتور، گزارش‌گیری و ارسال گزارش خارجی را یکپارچه می‌کند.

معماری مبتنی بر خط لوله، لایه‌بندی شفاف سرویس‌ها و استفاده از کش و صف،
سبب شده است که سامانه نه‌تنها از منظر توسعه‌پذیری و آزمون‌پذیری مناسب
باشد، بلکه از دید بهره‌برداری نیز قابلیت رشد داشته باشد. همچنین، پشتیبانی
از ارائه‌دهنده‌ای مانند \lr{Gemini} که از قرارداد بومی متفاوتی استفاده
می‌کند، نشان داد که معماری انتخاب‌شده برای تطبیق با ناهمگنی ارائه‌دهندگان
کفایت دارد. در ارزیابی انجام‌شده، 55 آزمون خودکار با 160
\lr{Assertion} با موفقیت اجرا شدند؛ بنابراین می‌توان گفت که نسخه‌ی فعلی
از نظر درستی مسیرهای اصلی، شواهد مناسبی ارائه می‌دهد. با این حال، این
نتیجه‌گیری به‌معنای اثبات کارایی میدانی در مقیاس تولید نیست.

\section{دستاوردها}

مهم‌ترین دستاورد این پروژه، طراحی و پیاده‌سازی یک درگاه سازگار با
\lr{OpenAI} برای دسترسی به چند ارائه‌دهنده‌ی هوش مصنوعی در قالب شش مسیر
اصلی هوش مصنوعی است. این دستاورد تنها به لایه‌ی درگاهی محدود نمی‌شود؛
زیرا در کنار آن سازوکاری یکپارچه برای مدیریت کاربران، مشتریان و کلیدهای
\lr{API} نیز ایجاد شده است.

از منظر مالی و عملیاتی نیز پروژه به یک مدل ترکیبی شامل کیف پول، اشتراک،
فاکتور و گزارش مصرف رسیده و منطق تخمین و اندازه‌گیری مصرف را هم برای
توکن و هم برای معیارهای غیرتوکنی توسعه داده است. افزوده شدن گزارش‌گیری
داخلی و ارسال لاگ ساخت‌یافته به سرویس‌های خارجی، در کنار ۲۸ فایل آزمون و
۵۵ آزمون خودکار، باعث شده سامانه فقط یک نمونه‌ی مفهومی نباشد. وجود
\lr{Playground} و مستندات \lr{API} نیز استفاده و توسعه‌ی بعدی آن را
آسان‌تر کرده است.

\section{محدودیت‌ها}

با وجود تحقق اهداف اصلی، نسخه‌ی فعلی سامانه همچنان دارای محدودیت‌هایی است:

ارسال \lr{OTP} در سطح منطق کامل شده است، اما هنوز به سرویس واقعی پیامک یا
ایمیل متصل نیست. به همین ترتیب، شارژ کیف پول نیز به درگاه پرداخت واقعی
متصل نشده و در نسخه‌ی فعلی به‌صورت دستی انجام می‌شود. قواعد راهبری
پیشرفته در جدول \lr{routing\_rules} پیش‌بینی شده‌اند، اما انتخاب
ارائه‌دهنده هنوز به‌صورت صریح توسط مشتری انجام می‌شود و سامانه تصمیم‌گیری
هوشمند را بر عهده نگرفته است.

از سوی دیگر، پشتیبانی از \lr{Gemini} در همه‌ی نقاط پایانی چندرسانه‌ای
تعمیم نیافته است و ارزیابی عملکرد این پایان‌نامه نیز بیشتر بر تحلیل
معماری و آزمون‌های نرم‌افزاری تکیه دارد. بنابراین بنچ‌مارک میدانی
گسترده، آزمون بار و تطبیق بلندمدت با صورت‌حساب واقعی ارائه‌دهندگان هنوز
خارج از دامنه‌ی نسخه‌ی فعلی قرار می‌گیرند.

\section{کارهای آتی}

برای توسعه‌ی بعدی سامانه، چند مسیر روشن وجود دارد. مهم‌ترین آن‌ها اتصال
\lr{OTP} به سرویس‌های واقعی پیامک یا ایمیل و یکپارچه‌سازی با درگاه پرداخت
برای شارژ خودکار کیف پول و تسویه‌ی اشتراک است. در سطح معماری نیز
پیاده‌سازی راهبردهای هوشمند انتخاب ارائه‌دهنده، مانند انتخاب
ارزان‌ترین، سریع‌ترین یا کم‌بارترین مسیر، می‌تواند گام بعدی طبیعی باشد.

افزودن قابلیت \lr{streaming} برای پاسخ‌های برخط، توسعه‌ی پشتیبانی
چندرسانه‌ای کامل‌تر برای همه‌ی ارائه‌دهندگان، اضافه کردن معیارهای
مشاهده‌پذیری پیشرفته‌تر مانند ردگیری توزیع‌شده و داشبوردهای عملیاتی، و
در نهایت انجام آزمون‌های بار و تحلیل تجربی عملکرد در سناریوهای واقعی
سازمانی، مسیرهایی هستند که می‌توانند نسخه‌ی فعلی را به یک سامانه‌ی
آماده‌تر برای بهره‌برداری گسترده تبدیل کنند.

در مجموع، پروژه‌ی حاضر نشان می‌دهد که طراحی یک سامانه‌ی واسط برای
تجمیع سرویس‌های هوش مصنوعی، نه‌فقط از منظر فنی شدنی است، بلکه می‌تواند
به‌عنوان مبنایی قابل توسعه برای مدیریت مصرف هوش مصنوعی در سطح سازمان
مورد استفاده قرار گیرد؛ مشروط بر آن‌که در گام بعدی، یکپارچه‌سازی‌های
بیرونی و ارزیابی میدانی آن تکمیل شود.
